% !TeX root = ../../main.tex
\subsection{個数と代金の問題}

\begin{techniquebox}[tech:02]{\textbf{個数と代金の問題}}
    \textbf{（単価）$\times$（個数）= \fitblankbf[代金]{}} という考え方を使う!
\end{techniquebox}

    基本は、Technique \ref{tech:02} の考え方を使って、\uwave{払う代金は変わらない}ことから、

    \[\textbf{(代金の合計) = (代金の合計)}\]

    \vspace{1.5em}
    という式を立てる。または、連立方程式を立てるときは、\uwave{購入する個数は変わらない}ことから、

    \[\textbf{(個数の合計) = (個数の合計)}\]

    という式を立てることも多い。

    \begin{example}[【例題 1】]
        同じノートを3冊と1本70円のペンを4本買うと、
        代金は940円になりました。ノート1冊の値段を求めなさい。
    \end{example}

    \begin{proof}\mbox{}\\
        \[（ノート3冊とペン4本の合計代金）=（支払う代金）\]

        という方針にしたがって、方程式を立てる。

        \begin{hiddenanswer}
            \begin{align*}
                \text{ノート1冊の値段を $x$ 円とすると、}\\
                3x + 70 \times 4 &= 940 \\
                3x + 280 &= 940 \\
                3x &= 660 \\
                x &= 220
            \end{align*}

            \finalanswer{答え：220円}
        \end{hiddenanswer}
    \end{proof}

\newpage
